\subsection{评分假设向资源配置传播的路径}
问题一评分变化会通过 $Q_0$ 改变质量成本的起算位置，而质量指数变化会改变目标函数本身，两类扰动不能混为同一种误差。三组权重网格得到的 $Q_0$ 范围约为 $0.455$--$0.511$，它是确定性评分方案的取值范围。问题三另外采用 $Q_0$ 上下 $20\%$ 的压力情景，覆盖范围更宽；该压力情景不表示基线质量存在已估计的 $20\%$ 统计误差。

表~\ref{tab:q3-sensitive-full} 除质量水平外，还列出中档预算的参数量、数据量和质量增量。只比较 $Q^*$ 容易忽略成本基线移动：相同 $Q^*$ 在不同 $Q_0$ 下，所需额外质量成本不同；而相近的质量水平也可以对应不同的规模与数据配置。基线提高时，部分原本需要支付的质量成本不再计入增量支出，节省的预算可以重新分配，不能据此推断现实中基线提高没有代价。

\begin{table}[!htbp]
\centering
\songti\zihao{-4}
\setlength{\tabcolsep}{4pt}
\caption{中档预算下的配置敏感性：绝对质量与增量质量}
\label{tab:q3-sensitive-full}
\begin{tabular}{lrrrr}
\toprule
情景 & $N^*$ (B) & $D^*$ (B) & $Q^*$ & $Q^*-Q_0$ \\
\midrule
基准 & 5.483 & 261.38 & 0.9665 & 0.4786 \\
$Q_0$ 降20\% & 5.499 & 260.04 & 0.9672 & 0.5769 \\
$Q_0$ 升20\% & 5.455 & 263.83 & 0.9651 & 0.3797 \\
$\theta=0.2$ & 5.127 & 293.60 & 0.8215 & 0.3336 \\
$\theta=0.8$ & 5.632 & 250.32 & 1.0000 & 0.5121 \\
\bottomrule
\end{tabular}
\end{table}


对质量指数的扰动，正确的比较对象是决策变化及其敏感程度。$\theta$ 越大，相同质量提升在目标函数中对应更强的数据项下降，候选质量通常更高；但不同 $\theta$ 下报告的损失来自不同函数，不能把二者的差值称为某个算法真实获得的提升。若要评价“按错误指数做决策”的损失代价，需要选定同一个参考损失函数，将所有情景配置回代该参考函数；这属于新的稳健决策实验，与本表的条件配置比较不同。

\subsection{复核层级与结论的证据边界}
本文将复核分成四个层级。第一层为输入与定义核对，包括字段解码、归一化参照、样本去留和单位换算。第二层为代数与实现核对，包括百分位总体均值恒等式、$Q=1$ 退化、解析偏导与有限差分一致，以及固定质量解析解与数值解一致。第三层为独立算法核对，包括加密网格、SLSQP 交叉验证及预算回代。第四层为外部预测核对，包括跨来源损失误差、分组留出、桥接留一和逐月前沿回测。

前三层通过，表明程序按照声明的公式运行且主要数值误差受到控制；只有第四层直接涉及模型在未参与拟合的观测上的表现。即使第四层通过，结论也仍限于相应测试分布与预测期限。例如在一个月后的回测表现良好，不足以证明两年后区间经过校准；同尺度配比排名较准，也不等于跨尺度绝对损失准确。区分这些层级，可以避免把“程序无错误”误写成“模型已经验证”。

对新增图表，分来源 RMSE 已按样本数复原总体误差；有限增量替代曲线已回代同一目标损失；固定质量解析公式已对照主求解器保存值；滚动预测的逐月误差已重新聚合并复原原有 MAE。所有这些复核都围绕明确的数值对应关系进行，没有因正文扩展而改变原始附件或主模型的既有计算结果。
